\section{Method}
\label{sec:method}

\begin{figure*}[t]
    \centering
    \includegraphics[width=\textwidth]{figures/full-pipe-ver7.png}
    \caption{Overview of \textsc{DeepTutor}. The system unifies two task pipelines through a Hybrid Personalization layer. \textbf{Static Knowledge Grounding} (center) constructs multi-modal knowledge graphs and vector databases from course materials, enabling mixed graph--embedding retrieval. \textbf{Dynamic Personal Memory} (bottom) maintains a trace forest accessible via a programmatic \emph{TraceToolkit}, and distills interaction histories into three evolving profile dimensions. \textbf{Personalized Problem Solving} (left) proceeds through Investigation~(\ding{172}), Step-by-step Solving~(\ding{173}), and Iterative Writing~(\ding{174}). \textbf{Personalized Question Generation} (right) proceeds via Idea Generation~(\ding{175}) and a Critic-Guided Q-A Pair Generation loop~(\ding{176}). Both pipelines record interaction traces that continuously refine the learner profile, forming a closed tutoring cycle.}
    \label{fig:overview}
\end{figure*}

DeepTutor is built around two sub-pipes with a common personalization engine, as shown in figure~\ref{fig:overview}.
The \textit{Hybrid Personalization Engine} (\S\ref{sec:personalization}) couples a static knowledge grounding---multi-modal knowledge graphs and vector databases built from course materials---with a dynamic memory that maintains an evolving learner profile.
Building upon this substrate, the \textit{problem-solving pipeline} (\S\ref{sec:solve}) produces citation-grounded explanations through three stages, while the \textit{question-generation pipeline} (\S\ref{sec:generate}) produces difficulty-calibrated practice items.
Each completed interaction feeds back into the learner profile, creating a closed tutoring cycle (\S\ref{sec:system}).

% ===========================================================================
\subsection{Preliminaries}
\label{sec:prelim}
% ===========================================================================

\paragraph{Task Formulation.}
DeepTutor addresses two complementary tutoring tasks:
(i)~\emph{problem solving}, producing a citation-grounded explanation $a$ for a learner question $q$, calibrated to the student's proficiency; and
(ii)~\emph{question generation}, producing customized practice items for a topic $\tau$ that target the learner's specific gaps.
Both tasks share the personalization substrate described below.

\paragraph{Notation.}
Course materials are indexed into a knowledge base $\mathcal{K}$ with dual representations for structured and semantic retrieval (\S\ref{sec:rag}).
The system maintains a learner profile $\mathcal{D} = (\mathcal{D}_s, \mathcal{D}_w, \mathcal{D}_r)$ capturing interaction summaries, diagnosed weaknesses, and tutor self-reflections, continuously updated by specialized memory agents (\S\ref{sec:memory}).
Before every agent step, a personalization context $\mathcal{C}_{\text{mem}}$ is assembled from $\mathcal{D}$ and the trace forest (the latter queried via a dedicated \emph{TraceToolkit}), yielding a compact, role-specific excerpt that steers the agent's behavior (\S\ref{sec:injection}).

Algorithm~\ref{alg:loop} gives a unified view of the full tutoring cycle.
Each session begins with profile injection and mixed retrieval, then branches into either problem solving (\ding{172}--\ding{174}) or question generation (\ding{175}--\ding{176}).
The resulting interaction trace is appended to the trace forest and used to update the learner profile.

\begin{algorithm}[t]
\caption{Closed-Loop Personalized Tutoring}\label{alg:loop}
\small
\begin{algorithmic}[1]
\Require learner profile $\mathcal{D}^{(0)}$;\; trace forest $\mathcal{F}^{(0)} = \emptyset$
\For{session $j = 1, 2, \ldots$}
    \State $\mathcal{C}_{\text{mem}} \gets \Call{ProfileInject}{\mathcal{D}^{(j{-}1)}, \mathcal{F}^{(j{-}1)}, \text{task}_j}$
    \State $\mathcal{C}_{\text{rag}} \gets \textsc{RagTool}(\text{task}_j)$
    \If{$\text{task}_j = (\textsc{Solve},\, q)$}
        \State $\mathcal{P}{=}\langle s_1,\dots,s_K\rangle \gets \Call{Plan}{q,\;\mathcal{C}_{\text{rag}},\;\mathcal{C}_{\text{mem}}}$ \Comment{\ding{172}}
        \For{$k = 1$ \textbf{to} $K$} \Comment{\ding{173}}
            \While{$s_k$ not resolved}
                \State $(\alpha_t,\, \mathbf{s}_t) \gets \Call{ToolAgent}{q, \mathcal{P}, \mathcal{S}, \mathcal{C}_{\text{mem}}}$
                \State $\mathcal{S} \gets \mathcal{S} \cup \{\mathbf{s}_t\}$
                \If{$\alpha_t = \textsc{Replan}$}
                    \State $\mathcal{P} \gets \Call{Revise}{\mathcal{P}, \mathcal{S}}$
                \EndIf
            \EndWhile
            \State $\mathcal{S} \gets \Call{Compress}{\mathcal{S}, k}$
        \EndFor
        \State $a \gets \Call{IterativeWrite}{\mathcal{S},\;\mathcal{C}_{\text{mem}}}$ \Comment{\ding{174}}
        \State $T_j \gets \Call{BuildTrace}{q, a}$
    \Else \Comment{$\text{task}_j = (\textsc{Generate},\, \tau)$}
        \State $\mathcal{I} \gets \Call{IdeaAgent}{\tau,\;\mathcal{C}_{\text{rag}},\;\mathcal{C}_{\text{mem}}}$ \Comment{\ding{175}}
        \State $\mathcal{T}_{1..N} \gets \Call{TopKSelect}{\mathcal{I}}$
        \For{$i = 1$ \textbf{to} $N$} \Comment{\ding{176}}
            \Repeat
                \State $(q_i, a_i) \gets \Call{Generator}{\mathcal{T}_i,\;\mathcal{C}_{\text{rag}},\;\mathcal{C}_{\text{mem}}}$
                \State $(\textit{pass}, f_i) \gets \Call{Validator}{q_i, a_i, \mathcal{T}_i}$
                \If{$\lnot\,\textit{pass}$} $\mathcal{T}_i \gets \Call{Revise}{\mathcal{T}_i, f_i}$
                \EndIf
            \Until{\textit{pass}}
        \EndFor
        \State $T_j \gets \Call{BuildTrace}{\tau, \{(q_i, a_i)\}}$
    \EndIf
    \State $\mathcal{F}^{(j)} \gets \mathcal{F}^{(j{-}1)} \cup \{T_j\}$
    \State $\mathcal{D}^{(j)} \gets \Call{MemoryUpdate}{\mathcal{D}^{(j{-}1)}, T_j}$ 
\EndFor
\end{algorithmic}
\end{algorithm}

% ===========================================================================
\subsection{Hybrid Personalization Engine}
\label{sec:personalization}
% ===========================================================================

Both the solver and the generator require two types of knowledge: \emph{domain expertise} and \emph{learner awareness}. DeepTutor provides the former through Static Knowledge Grounding and the latter through Dynamic Memory, unified by a Profile Injection mechanism that delivers tailored context to each specific stage.

\subsubsection{Static Knowledge Grounding}
\label{sec:rag}

Real-world course materials are inherently multimodal: textbooks interleave prose with figures, equations, and data tables, each encoding domain knowledge that is difficult to faithfully flatten into plain text.
To preserve this richness, we perform multimodal knowledge unification, decomposing source documents into atomic content units $\{c_j=(t_j, x_j)\}$, where $t_j\!\in\!\{\text{text}, \text{image}, \text{table}, \text{equation}\}$ denotes the modality and $x_j$ the raw payload~\cite{wang2024mineruopensourcesolutionprecise}.
Dedicated extractors handle each modality so as to retain its internal structure: figures keep their captions, equations stay anchored to surrounding definitions, and tables preserve header--cell relationships.

The resulting units are indexed through a dual-graph construction: a cross-modal KG grounds non-textual units within their contextual neighborhood via entity and relation extraction, while a text-based KG captures fine-grained semantic relationships from prose~\cite{guo2025raganythingallinoneragframework}.
The two graphs are merged through entity alignment into a unified knowledge graph $\mathcal{G}=(\mathcal{V},\mathcal{E})$ that encodes both cross-modal associations and textual semantics.
In parallel, dense encoders map all content units into modality-specific vector stores (a text VDB and a multi-modal VDB), unified into a combined embedding index $\mathcal{B}$.
Together, $\mathcal{G}$ and $\mathcal{B}$ form the complete retrieval index over the knowledge base $\mathcal{K}$.

At query time, we introduce cross-modal hybrid retrieval that operates over both indices in parallel: graph-based retrieval traverses $\mathcal{G}$ to capture explicit relationships and multi-hop reasoning paths, while embedding-based retrieval searches $\mathcal{B}$ to surface semantically relevant content beyond the graph's structural coverage.
The resulting candidate sets are merged via reciprocal rank fusion, scoring each passage $d$ as $\textstyle\sum_{r}{1}/({k+\mathrm{rank}_r(d)})$, where $k$ is a smoothing constant~\citep{cormack2009reciprocal}.
Retrieved passages are deduplicated and truncated to a context budget before injection as domain grounding $\mathcal{C}_{\text{rag}}$.

\subsubsection{Dynamic Personal Memory}
\label{sec:memory}

While the knowledge base captures \emph{what the course materials contain}, the dynamic personal memory captures \emph{how the individual learner has engaged with them}.
We further propose \textit{Trace Forest}, a structured memory $\mathcal{F} = \{T_1, T_2, \ldots\}$ in which each tree $T_j = (V_j, E_j)$ records one complete tutoring interaction as a multi-resolution, semantically searchable artifact.
Every node $v \in V_j$ is a triple $(\texttt{type}, \texttt{content}, \mathbf{e}_v)$ where \texttt{type} identifies the granularity level, \texttt{content} stores the payload, and $\mathbf{e}_v$ is a dense embedding, forming a three-level hierarchy. \textit{Level~1} captures the session-level input and a global summary; \textit{Level~2} records intermediate planning units produced by task decomposition; and \textit{Level~3} stores fine-grained execution records including tool invocations, evidence, and validation outcomes.
The concrete node types at each level are task-specific and detailed alongside the respective pipelines (\S\ref{sec:solve}, \S\ref{sec:generate}).

Rather than treating the forest as a passive data store, we expose it to all pipeline agents through a programmatic \emph{TraceToolkit} with three operations: \textsc{SearchTrace} embeds a query and retrieves the most relevant nodes across $\mathcal{F}$ via cosine similarity-based ANN search using a shared encoder; \textsc{ListTraces} enumerates traces filtered by time range, task type, or topic; and \textsc{ReadNodes} fetches the full content of specific nodes along with their ancestry path.
This interface lets each agent actively explore the learner's history with their specific demands, from coarse session summaries at Level~1 to fine-grained reasoning steps at Level~3.

\paragraph{Profile Construction.}
\label{sec:injection}
Raw traces must be interpreted into actionable tutoring signals.
We introduce three specialized memory agents to process each new trace tree $T_j$ in parallel, incrementally updating a dedicated dimension of the learner profile $\mathcal{D}$:
\begin{itemize}[nosep,leftmargin=*]
\item Session History $ \mathcal{D}_s$: distills topics covered and performance trends. This will record what has happened, providing general session memories.
\item User Weakness $\mathcal{D}_w$: compares learner responses against reference solutions, labeling each weakness as active or resolved to maintain a prioritized gap inventory, emphasizing where to focus.
\item Self-Reflection $\mathcal{D}_r$: critiques the system's own prior outputs for pedagogical alignment, producing actionable improvement notes, indicating how to improve.
\end{itemize}
Each agent receives the new trace along with its current profile dimension and produces an updated version through incremental revision. To prevent error accumulation, this process employs a \textit{confidence-gated update} mechanism, where state transitions (e.g., marking a weakness as resolved) require explicit evidence grounding, preserving long-term trends while incorporating recent signals.

Before each agent step, the system assembles a personalization context $\mathcal{C}_{\text{mem}}$ through two channels:
(i)~\emph{active trace retrieval} via the TraceToolkit, retrieving the top-$k$ most relevant nodes with $k$ allocated proportionally across levels to balance broad session context with fine-grained precedents;
(ii)~\emph{profile excerpting}, where role-specific slices of $\mathcal{D}$ are drawn based on the requesting agent's function to decouple their attention (e.g., the planner receives $\mathcal{D}_s$ and $\mathcal{D}_w$; the writer receives $\mathcal{D}_r$; generation agents receive $\mathcal{D}_w$ alongside historical question patterns). The total token budget is dynamically partitioned between $\mathcal{C}_{\text{rag}}$ and $\mathcal{C}_{\text{mem}}$ (see Appendix for a worked example).

% ===========================================================================
\subsection{Personalized Problem-Solving}
\label{sec:solve}
% ===========================================================================

Prior tool-augmented agents~\citep{yao2023reactsynergizingreasoningacting, chen2026iterresearch} pack planning, retrieval, and exposition into a full reasoning loop.
This suffices for short factoid queries but becomes a bottleneck in tutoring, where thorough investigation and personalized presentation compete for the same context budget.
DeepTutor disentangles them into three incremental stages (\ding{172}--\ding{174} in Algorithm~\ref{alg:loop}), guided by \textit{scaffolding theory} to provide step-by-step cognitive support.

\paragraph{Stage \ding{172}: Personalized Investigation.}
The planner decomposes $q$ into several meta-questions and launches a multi-source investigation: it retrieves relevant passages from $\mathcal{K}$ via the RAG tool, consults the trace forest via the trace toolkit for precedents from similar past interactions, and invokes auxiliary toolkits as needed.
Jointly conditioning on the assembled domain context $\mathcal{C}_{\text{rag}}$ and the learner context $\mathcal{C}_{\text{mem}}$, the planner produces a solving plan $\mathcal{P} = \langle s_1, \ldots, s_K \rangle$ of concrete sub-goals, each annotated with its step-level goal and a brief description.
A shared scratchpad $\mathcal{S}$ accumulates all subsequent reasoning artifacts.
Investigating before planning prevents sub-goals that are too vague or redundant with what the learner already understands (e.g., generating ``review chain rule in trigonometry'' rather than a vague ``review calculus'').

\paragraph{Stage \ding{173}: Step-by-step Problem Solving.}
Each sub-goal $s_k$ is executed through a ReAct-style loop, ~\citep{yao2023reactsynergizingreasoningacting}, conditioned on both the domain grounding $\mathcal{C}_{\text{rag}}$ and the learner context $\mathcal{C}_{\text{mem}}$.
At every iteration the tool agent produces a reasoning trace $r_t$, selects an action $\alpha_t$ from a shared tool suite $\mathcal{A}$ (encompassing knowledge-base retrieval, trace-forest lookup, sandboxed code execution, web search, and deep reasoning), observes the output $o_t$, and records a \emph{self-note} $n_t$.
Unlike the reasoning trace, which captures the full deliberation process, the self-note distills the step's outcome into a concise takeaway (e.g., ``retrieved formula confirmed'') that subsequent steps can reference without re-reading verbose intermediates.
When the current plan proves inadequate, the agent triggers \emph{adaptive replanning}, revising remaining sub-goals while preserving completed work.

Hierarchical compression (line 14 in Algorithm~\ref{alg:loop}) prevents the scratchpad from exceeding the model's context window.
Once sub-goal $s_k$ is resolved, an LLM-based summarizer distills its full trace into a compact digest that retains key conclusions and evidence pointers while discarding verbose intermediate steps.
As the plan advances, earlier sub-goals are represented by increasingly concise digests, freeing context capacity for deeper reasoning on later sub-goals.

\paragraph{Stage \ding{174}: Evidence-based Iterative Writing.}
Rather than concatenating raw tool outputs, the writing stage extracts structured evidence bundles from the scratchpad and constructs the answer through successive draft-refine iterations, reconciling potentially conflicting findings.
The learner context $\mathcal{C}_{\text{mem}}$
% ---including diagnosed weaknesses from $\mathcal{D}_w$ and the system's pedagogical self-reflections from $\mathcal{D}_r$---
steers both the depth and the tone of the explanation within the learner's \textit{Zone of Proximal Development}: \textit{beginners} receive step-by-step derivations with scaffolding, while \textit{proficient learners} receive focused summaries that foreground key insights.
Every claim is grounded with traceable citations back to $\mathcal{K}$.

% ===========================================================================
\subsection{Personalized Question Generation}
\label{sec:generate}
% ===========================================================================

A tutor must also pose the right questions: ones that are meaningful, factually grounded, and calibrated to the individual learner.
We propose a two-stage architecture (\ding{175}--\ding{176} in Algorithm~\ref{alg:loop}) that separates \emph{what to ask} from \emph{how to ask and verify}.

\paragraph{Stage \ding{175}: Personalized Idea Generation.}
Rather than generating questions directly from a bare topic, DeepTutor first maps the conceptual landscape around $\tau$ through the lens of the individual learner.
The idea agent jointly conditions on domain knowledge $\mathcal{C}_{\text{rag}}$ and learner context $\mathcal{C}_{\text{mem}}$---including past mistakes, difficulty trends, and related problems surfaced via the TraceToolkit---to synthesize a diverse pool of candidate ideas, each characterized by a target concept, a question format (MCQ, written-response, or coding), and a personalized rationale grounded in the learner's diagnosed gaps.
A two-stage selection then filters ill-formed or redundant candidates and ranks survivors along quality dimensions such as clarity, relevance, and diversity, yielding structured question templates $\{\mathcal{T}_i\}$ that serve as blueprints for the subsequent generation stage.

\paragraph{Stage \ding{176}: Critic-Guided Q-A Pair Generation.}
For each template $\mathcal{T}_i$, a \textbf{generator} produces a complete question--answer--explanation triple $(q_i, a_i, e_i)$, drawing on $\mathcal{C}_{\text{rag}}$ for factual grounding and $\mathcal{C}_{\text{mem}}$ to calibrate difficulty and phrasing to the learner's level.
A structurally separated \textbf{validator} then challenges the output through two-level quality assurance: an LLM-based reasoning check evaluates template alignment, factual correctness, and pedagogical soundness, while computational questions undergo additional verification via sandboxed code execution.
This dual-agent separation is critical: because the validator shares no reasoning chain with the generator, it must independently verify correctness, catching errors that self-evaluation would miss.
Failed pairs receive structured diagnostic feedback $f_i$; the template is revised and the pair regenerated until both pedagogical and factual constraints are satisfied.

% ===========================================================================
\subsection{Closed-Loop Tutoring Cycle}
\label{sec:system}
% ===========================================================================

As shown in Algorithm~\ref{alg:loop}, the components above form a self-reinforcing cycle.
After each task, whether a problem-solving session or a question-generation episode, the pipeline records a new trace tree $T_j$ into the trace forest $\mathcal{F}$, and the three memory agents update the learner profile in parallel, producing $\mathcal{D}^{(j)}$ that immediately informs subsequent interactions.
The key property is \textbf{bidirectional task coupling}, which operationalizes \textit{formative assessment}: weaknesses diagnosed during problem solving propagate to $\mathcal{D}_w$ and directly shape which questions are generated next; conversely, the learner's performance on generated questions refines $\mathcal{D}_s$ and $\mathcal{D}_r$, improving future explanations.
Because every agent can actively query the trace forest through the TraceToolkit, this feedback is not limited to profile-level summaries: agents can retrieve fine-grained precedents from any prior interaction, enabling increasingly nuanced personalization as the interaction history grows.
